%%%%%%%%%%%%%%%%%%%%%%%%%%%%%%%%%%%%%%%%%%%%%%%%%%%%%%%%%%%%%%%%%%%%%%
% Lớp 12 | Phương trình mặt phẳng, đường thẳng, mặt cầu | Phương trình đường thẳng trong không gian
% Dạng 2H5?2-1: Câu hỏi lý thuyết
% Số câu: 15  (sinh tự động bởi tach_cau.py; từ nay thầy sửa trực tiếp file này)
%%%%%%%%%%%%%%%%%%%%%%%%%%%%%%%%%%%%%%%%%%%%%%%%%%%%%%%%%%%%%%%%%%%%%%

%%%%%%%%%%%%%%%% TRẮC NGHIỆM 4 PHƯƠNG ÁN %%%%%%%%%%%%%%%%

% nguồn: Phan1_Lop_12_Hinh_hoc_Chuong_5
\begin{ex}%[2H5N2-1]%[ma:e732d788]
	Trong không gian $Oxyz$, cho đường thẳng $\Delta \colon \dfrac{x-1}{3}=\dfrac{y+2}{-5}=\dfrac{z-4}{1}$. Trong các điểm sau đây, điểm nào \textbf{không} thuộc  đường thẳng $\Delta$?
	 \choice
		{\True  $M(4;-7;-2)$}
		 {$N(7;-12;6)$}
		 {$P(10;-17;7)$}
		 {$Q(13;-22;8)$}
	 \loigiai{
		Thế tọa độ $M(4;-7;-2)$ vào phương trình đường  thẳng $\Delta$ ta được\\
		$\dfrac{4-1}{3}=\dfrac{-7+2}{-5}=\dfrac{-2-4}{1} \Leftrightarrow 1=1=-6$ (vô  lý).
		Suy ra $M$ không thuộc đường thẳng  $\Delta$.\\
		Thế tọa độ  $N(7;-12; 6)$ vào phương trình đường thẳng $\Delta$ ta được\\
		 $\dfrac{7-1}{3}=\dfrac{-12+2}{-5}=\dfrac{6-4}{1} \Leftrightarrow 2=2=2$ (đúng).
		Suy ra $N$ thuộc  đường thẳng $\Delta$.\\
		Thế tọa độ $P(10;-17;7)$ vào phương  trình đường thẳng $\Delta$ ta được\\ $\dfrac{10-1}{3}=\dfrac{-17+2}{-5}=\dfrac{7-4}{1} \Leftrightarrow 3=3=3$ (đúng).
		Suy ra $P$ thuộc đường thẳng  $\Delta$.\\
		Thế tọa độ $Q(13;-22;8)$ vào phương trình đường  thẳng $\Delta$ ta được\\ $\dfrac{13-1}{3}=\dfrac{-22+2}{-5}=\dfrac{8-4}{1} \Leftrightarrow 4=4=4$ (đúng).
	Suy ra $Q$ thuộc đường  thẳng $\Delta$.}
\end{ex}

% nguồn: Phan1_Lop_12_Hinh_hoc_Chuong_5
\begin{ex}%[2H5N2-1]%[ma:bc1ee9c8]
	Trong không gian $Oxyz$, điểm nào dưới đây thuộc đường thẳng  $\heva{&x=1-t\\&y=5+t\\&z=2+3t}$
	\choice{$Q(-1;1;3)$} 
		{$N(1;-1;3)$} 
		{$A(-2;5;-2)$} 
		{\True $P(1;5;2)$} 
	\loigiai{Điểm thuộc đường thẳng  $\heva{&x=1-t\\&y=5+t\\&z=2+3t}$ là $P(1;5;2)$.}
\end{ex}

% nguồn: Phan1_Lop_12_Hinh_hoc_Chuong_5
\begin{ex}%[2H5N2-1]%[ma:514b059c]
	Trong không gian với hệ toạ độ $Oxyz$, phương trình nào sau đây là  phương trình chính tắc của đường thẳng?
	\choice 
		{\True  $\dfrac{x-9}{3}=\dfrac{y-1}{11}=\dfrac{z+7}{14}$}
		{$\dfrac{x-1}{4y}=\dfrac{y-2}{5}=\dfrac{z+3}{6}$} 
		 {$\dfrac{x+7}{10}=\dfrac{y+8}{11z}=\dfrac{z+9}{12}$}
		{$\dfrac{x-13}{16}=\dfrac{y-14}{17}=\dfrac{z-15}{18z}$} 
	\loigiai{ 
		Phương trình chính tắc của đường thẳng  là $\dfrac{x-9}{3}=\dfrac{y-1}{11}=\dfrac{z+7}{14}$.
	} 
\end{ex}

% nguồn: Phan1_Lop_12_Hinh_hoc_Chuong_5
\begin{ex}%[2H5N2-1]%[ma:88ef5b5b]
	Trong  không gian $Oxyz$, đường thẳng $\heva{&x=1+2t\\&y=3-t\\&z=1-t}$ \textbf{không} đi qua điểm nào dưới đây?
	 \choice{$Q(3;2;0)$}
		 {$N(5;1;-1)$}
		{$A(1;3;1)$} 
		{\True  $P(0;1;2)$}
	\loigiai{ 
		Điểm không thuộc đường  thẳng $\heva{&x=1+2t\\&y=3-t\\&z=1-t}$ là $P(0;1;2)$.
	} 
\end{ex}

% nguồn: Phan1_Lop_12_Hinh_hoc_Chuong_5
\begin{ex}%[2H5N2-1]%[ma:ae411654]
	Trong không gian với hệ tọa độ $Oxyz$, phương trình nào sau đây có  dạng là phương trình tham số của đường thẳng?
	 \renewcommand{\dotEX}{}
	\choice 
		{\True $\heva{&x=1+2t\\&y=-3+4t\\&z=5-7t.}$} 
		 {$\heva{&x=3+t^2\\&y=9+2t\\&z=-6+11t.}$}
		{$\heva{&x=9+17t\\&y=8+2t\\&z=-12+t^2.}$} 
		{$\heva{&x=15+t\\&y=25+\sqrt{t}\\&z=35-t.}$} 
	 \loigiai{
	Phương trình tham số của đường thẳng có dạng  $\heva{&x=x_0+at\\&y=y_0+bt\\&z=z_0+ct} (\text{$t$}$ là tham số).}
\end{ex}

% nguồn: Phan1_Lop_12_Hinh_hoc_Chuong_5
\begin{ex}%[2H5N2-1]%[ma:f05e97ee]
	Trong không gian $Oxyz$, đường thẳng $d\colon \dfrac{x-1}{2}=\dfrac{y+1}{-1}=\dfrac{z}{2}$ đi qua điểm  nào dưới đây?
	 \choice
		 {\True $(1;-1;0)$}
		 {$(1;0;1)$}
		{$(1;-1;1)$} 
		{$(0;0;1)$} 
	 \loigiai{
		Đường thẳng $d\colon \dfrac{x-1}{2}=\dfrac{y+1}{-1}=\dfrac{z}{2}$ đi qua điểm có  tọa độ $(1;-1;0)$.
	 }
\end{ex}

% nguồn: Phan1_Lop_12_Hinh_hoc_Chuong_5
\begin{ex}%[2H5N2-1]%[ma:7ca98cc0]
	Trong không gian $Oxyz$, cho đường thẳng $d \colon \dfrac{x-1}{2}=\dfrac{y-2}{3}=\dfrac{z+1}{-1}$. Điểm nào  sau đây thuộc $d$?
	\choice{$Q(-1;-2;1)$} 
		 {$N(-1;3;2)$}
		 {$A(2;3;-1)$}
		{\True $P(1;2;-1)$} 
	\loigiai{ 
	Điểm thuộc đường thẳng $d \colon \dfrac{x-1}{2}=\dfrac{y-2}{3}=\dfrac{z+1}{-1}$  là $P(1;2;-1)$.}
\end{ex}

% nguồn: Phan1_Lop_12_Hinh_hoc_Chuong_5
\begin{ex}%[2H5N2-1]%[ma:db385d1e]
	Trong không gian $Oxyz$, điểm nào dưới  đây thuộc đường thẳng $d \colon \dfrac{x+1}{-1}=\dfrac{y-2}{3}=\dfrac{z-1}{3}$?
	 \choice
		{$A(1;  2; 1)$}
		 {$N(-1; 3; 2)$}
		{\True $P(-1;  2; 1)$}
		{$Q(1;-2;-1)$} 
	\loigiai{Điểm $P(-1; 2; 1)$ thuộc đường thẳng $d \colon \dfrac{x+1}{-1}=\dfrac{y-2}{3}=\dfrac{z-1}{3}$.} 
\end{ex}

% nguồn: Phan1_Lop_12_Hinh_hoc_Chuong_5
\begin{ex}%[2H5N2-1]%[ma:3d728806]
	Trong không gian $Oxyz$, điểm nào dưới đây thuộc đường thẳng $d$  có phương trình $\dfrac{x-1}{-1} = \dfrac{y-2}{3} = \dfrac{z+1}{3}$?
	 \choice
		 {$N(-1;4;3)$}
		 {\True $P(1;2;-1)$}
		 {$(1;2;1)$}
		{$Q(1;-2;2)$} 
	 \loigiai{Ta có
		\[\dfrac{x-1}{-1} = \dfrac{y-2}{3} = \dfrac{z+1}{3} \Leftrightarrow \dfrac{x-1}{-1}  = \dfrac{y-2}{3} = \dfrac{z-(-1)}{3}.\]
	 Vậy đường thẳng $d$ đi qua điểm $P(1;2;-1)$.}
\end{ex}

% nguồn: Phan1_Lop_12_Hinh_hoc_Chuong_5
\begin{ex}%[2H5N2-1]%[ma:f80a4954]
	Trong không gian $Oxyz$,  cho đường thẳng $d\colon\heva{&x=2- t \\ &y = -1+2t \\ &z = 3-3t}$. Điểm nào dưới đây thuộc $d$?
	 \choice
		{$Q(1; -2;  3)$}
		{$P(-1;  2; -3)$}
		{\True $M(2; -1;  3)$}
		 {$N(-2; 1; -3)$}
	\loigiai{Nhận ra, điểm  $M(2; -1; 3) \in d$ vì $\heva{&2=2- t \\ &-1 = -1+2t \\ &3 = 3-3t}\Rightarrow t = 0$.
	} 
\end{ex}


%%%%%%%%%%%%%%%% ĐÚNG/SAI %%%%%%%%%%%%%%%%

% nguồn: Phan1_Lop_12_Hinh_hoc_Chuong_5
\begin{ex}%[2H5H2-1]%[ma:4cc1c39d]
	Trong không gian $Oxyz$, cho hai đường thẳng $\Delta \colon \dfrac{x+1}{1}=\dfrac{y+5}{2}=\dfrac{z+9}{3}$, $\Delta'\colon  \dfrac{x+6}{3}=\dfrac{y+3}{2}=\dfrac{z}{1}$.
	 \choiceTF
		{\True Hai vectơ $\overrightarrow{u}(  1;2;3)$, $\overrightarrow{u'}( 3;2;1)$ lần lượt là vectơ chỉ phương của $\Delta $ và $\Delta' $}
		{Điểm $M( -1;-5;-9)$ không thuộc đường thẳng $\Delta $, điểm $M'( -6;-3;0)$ không  thuộc đường thẳng $\Delta' $}
		{\True $[\vec{u},\vec{u}']=(-4;8;-4)$} 
		{Hai đường thẳng $\Delta $ và $\Delta' $ chéo nhau} 
	 \loigiai{
		\begin{itemchoice}
			\itemch \textbf{Đúng}.\\Hai  vectơ $\overrightarrow{u}(1;2;3)$, $\overrightarrow{u'}( 3;2;1)$ lần lượt là vectơ chỉ phương của $\Delta$, ${\Delta}'$.
			 \itemch \textbf{Sai}.\\Điểm $M(-1;-5;-9)$ thuộc đường thẳng $\Delta$, điểm $M'( -6;-3;0)$ thuộc đường thẳng $\Delta'$.
			\itemch \textbf{Đúng}.\\Tích có hướng của hai vectơ  là $[\vec{u},\vec{u}']=(-4;8;-4)$.
			\itemch \textbf{Sai}.\\Gọi $M(-1;-5;-9)$ thuộc đường thẳng  $\Delta$, $M'(-6;-3;0)$ thuộc đường thẳng $\Delta'$.\\
			Khi đó, $[\overrightarrow{u},\overrightarrow{u'}]\cdot \overrightarrow{MM'}=0$. \\ 
			Vậy hai vectơ $\overrightarrow{u}$ và $\overrightarrow{u'}$ cắt  nhau.
	\end{itemchoice}}
\end{ex}


%%%%%%%%%%%%%%%% TRẢ LỜI NGẮN %%%%%%%%%%%%%%%%

% nguồn: Phan2_Lop_12_Hinh_hoc_Chuong_5
\begin{bt}%[2H5N2-1]%[ma:d07482ed]
	Trong không gian $Oxyz$, cho đường thẳng $\heva{&x=1+2t\\&y=2-t\\&z=3+2t}$ và các  điểm $M(-1;3;1)$, $N(-2;1;-2)$, $P(1;2;3)$, $Q(2;-1;2)$, $S(1;2;-3)$. Có mấy điểm \textbf{không} thuộc đường thẳng $d$ cho trước?
	 \shortans{3}
	 \loigiai{
		Thay tọa độ các điểm vào phương  trình đường thẳng $d$ ta thấy \\
		Các điểm $N(-2;1;-2)$, $Q(2;-1;2)$, $S(1;2;-3) \notin  d$.\\
	Vậy có $3$ điểm không thuộc $d$. } 
\end{bt}

% nguồn: Phan2_Lop_12_Hinh_hoc_Chuong_5
\begin{bt}%[2H5N2-1]%[ma:883ddecb]
	Trong không gian $Oxyz$, cho các điểm $M(-1;3;1)$, $N(-2;1;-2)$, $P(1;2;3)$, $Q(2;-1;2)$, $S(1;2;-3)$. Đường thẳng $d \colon \dfrac{x-1}{2}=\dfrac{y-2}{-1}=\dfrac{z-3}{2}$ đi qua bao nhiêu điểm trong số các điểm đã  cho?
	 \shortans{2}
	\loigiai{ 
		Thay tọa độ các điểm vào phương trình đường thẳng $d$ ta thấy $M(-1;3;1)$,  $P(1;2;3) \in d$. \\
	Vậy đường thẳng $d$ đi  qua $2$ điểm.}
\end{bt}

% nguồn: Phan2_Lop_12_Hinh_hoc_Chuong_5
\begin{bt}%[2H5H2-1]%[ma:95875241]
	Cho hình hộp $ABCD.A'B'C'D'$. Có bao nhiêu vectơ chỉ phương của đường thẳng $A B$ có  điểm đầu và điểm cuối là đỉnh của hình hộp?
	 \shortans{8}
	 \loigiai{
		Đường thẳng $\overrightarrow{A B}$ nhận các vectơ $\overrightarrow{A B}$, $\overrightarrow{B A}$, $\overrightarrow{A' B'}$, $\overrightarrow{B' A'}$, $\overrightarrow{C D}$, $\overrightarrow{D  C}$, $\overrightarrow{C' D'}$, $\overrightarrow{D' C'}$ là các vectơ chỉ phương.
	 Vậy có 8 vectơ thỏa mãn.}
\end{bt}


%%%%%%%%%%%%%%%% TỰ LUẬN %%%%%%%%%%%%%%%%

% nguồn: Phan2_Lop_12_Hinh_hoc_Chuong_5
\begin{bt}%[2H5N2-1]%[ma:110f42e9]
	Trong không  gian $Oxyz$, cho đường thẳng $\Delta \colon \heva{&x=2+3t\\&y=-1+t\\&z=5+-2t}$. Trong các điểm sau đây, điểm nào thuộc đường thẳng $\Delta$?
	\begin{listEX}[4]
		 \item $M(11;2;4)$
		\item  $N(5;0;3)$
		\item  $P(10;3;-3)$
		\item  $Q(8;-3;1)$.
	\end{listEX}
	 \loigiai{
		\begin{listEX}[1]
			\item Thế tọa độ $M(11;2;4)$ vào phương trình đường thẳng $\Delta$ ta  được:
			$\heva{&11=2+3t\\&2=-1+t\\&4=5-2t}\Leftrightarrow\heva{&t=3\\&t=3\\&t=\dfrac{1}{2}}$.\\ 
			Vì hệ  phương trình trên vô nghiệm nên $M$ không thuộc đường thẳng $\Delta$.\\
			\item Thế tọa độ  $N(5;0;3)$ vào phương trình đường thẳng $\Delta$ ta được: $\heva{&5=2+3t\\&0=-1+t\\&3=5-2t}\Leftrightarrow\heva{&t=1\\&t=1\\&t=1}$.\\
			 Vì hệ phương trình trên có nghiệm nên $N$ thuộc đường thẳng $\Delta$.\\
			\item Thế tọa độ $P(10;3;-3)$  vào phương trình đường thẳng $\Delta$ ta được:\\ $\heva{&10=2+3t\\&3=-1+t\\&-3=5-2t}\Leftrightarrow\heva{&t=\dfrac{8}{3}\\&t=4\\&t=4}$.\\
			Vì hệ phương trình trên vô nghiệm nên $P$  không thuộc đường thẳng $\Delta$.\\
			\item Thế tọa độ $Q(8;-3;1)$  vào phương trình đường thẳng $\Delta$ ta được:\\ $\heva{&8=2+3t\\&-3=-1+t\\&1=5-2t}\Leftrightarrow\heva{&t=2\\&t=-2\\&t=2}$.\\
			Vì hệ phương  trình trên vô nghiệm nên $Q$ không thuộc đường thẳng $\Delta$.
		\end{listEX}
	 }
\end{bt}
